% =====================================================
% 题型：正弦定理与余弦定理解答题 (q_triangle_cosine_rule_cos.tex)
% =====================================================
% =====================================================
% 难度：★★ (中档)
% =====================================================
\newcommand{\qTriangleCosineRuleCosStem}{在 \(\triangle ABC\) 中，\(a,b,c\) 分别为内角 \(A,B,C\) 所对的边，且满足 \(a\cos B + b\cos A = 2c\cos C\)。\\
（1）求 \(\cos C\)；\\
（2）若 \(c=3\)，且 \(\triangle ABC\) 的面积为 \(\dfrac{3\sqrt3}{4}\)，求 \(a,b\) 的值。}

\newcommand{\qTriangleCosineRuleCosAnswer}{（1）\(\cos C=\dfrac12\)；（2）\(a,b\) 为 \(1,2\)（次序可交换）。}

\newcommand{\qTriangleCosineRuleCosSolution}{%
  （1）由正弦定理
  \[
    \frac{a}{\sin A}
    = \frac{b}{\sin B}
    = \frac{c}{\sin C}
    = 2R,
  \]
  可令
  \[
    a=2R\sin A,
    \qquad
    b=2R\sin B,
    \qquad
    c=2R\sin C.
  \]
  代入原式得
  \[
    2R\sin A\cos B
    +2R\sin B\cos A
    =4R\sin C\cos C.
  \]
  两边同除以 \(2R\)，得
  \[
    \sin A\cos B+\sin B\cos A
    =2\sin C\cos C.
  \]
  即
  \[
    \sin(A+B)=2\sin C\cos C.
  \]
  因为 \(A+B=\pi-C\)，所以
  \[
    \sin(A+B)=\sin C.
  \]
  从而
  \[
    \sin C=2\sin C\cos C.
  \]
  又因为三角形内角满足 \(C\in(0,\pi)\)，故 \(\sin C>0\)，所以可约去 \(\sin C\)，得
  \[
    1=2\cos C.
  \]
  因此
  \[
    \cos C=\frac12,
    \qquad
    C=\frac{\pi}{3}.
  \]

  （2）由面积公式
  \[
    S_{\triangle ABC}
    =\frac12ab\sin C,
  \]
  得
  \[
    \frac12ab\cdot\frac{\sqrt3}{2}
    =\frac{3\sqrt3}{4}.
  \]
  所以
  \[
    ab=3.
  \]

  又由余弦定理
  \[
    c^2=a^2+b^2-2ab\cos C.
  \]
  代入 \(c=3\)、\(ab=3\)、\(\cos C=\dfrac12\)，得
  \[
    9=a^2+b^2-3.
  \]
  因此
  \[
    a^2+b^2=12.
  \]
  所以
  \[
    (a+b)^2=a^2+b^2+2ab=12+6=18,
  \]
  即
  \[
    a+b=3\sqrt2.
  \]
  此时 \(a,b\) 是方程
  \[
    x^2-3\sqrt2\,x+3=0
  \]
  的两个正根。
  解得
  \[
    x=\frac{3\sqrt2\pm\sqrt6}{2}.
  \]
  因此
  \[
    \{a,b\}
    =
    \left\{
      \frac{3\sqrt2+\sqrt6}{2},
      \frac{3\sqrt2-\sqrt6}{2}
    \right\}.
  \]
}

\newcommand{\qTriangleCosineRuleCosDiagnosis}{%
  （留白待收集学生作答后补充）%
}

\newcommand{\qTriangleCosineRuleCosTeachingNotes}{%
  第（1）问先用正弦定理统一成角，再用 \(\sin(A+B)=\sin C\) 化简；第（2）问用面积公式得到 \(ab\)，再用余弦定理得到 \(a^2+b^2\)，属于典型的“乘积 + 平方和”收束。%
}
